The unit of observation is the city-by-genre-by-year cell. The outcome is an indicator for whether
that cell emerges in a given year. The headline regressors are local spawning flow, target-genre
active bands, and target-genre multi-band musicians. The preferred specification also includes
broader city-level controls for overall active bands, nontrivial communities, and total multi-band
musicians. All predictors are transformed as one-year-lagged rolling averages and z-scored within
the estimation sample. Table \ref{tab:main_variables} reports the main variables used in the
current draft. Appendix Table \ref{tab:appendix_variable_crosswalk} reports the fuller variable
crosswalk, including the code-side variable names, units of observation, and the level at which
the rolling transformations are constructed before the full panel merge. This standardization is
useful later in the results section because the coefficients can be compared on a common
within-sample scale: each coefficient is the change in emergence probability associated with a
one-standard-deviation increase in the lagged predictor.

\begin{table}[htbp]
\centering
\caption{Main variables in the exact-year scene specification}
\label{tab:main_variables}
\small
\begin{tabularx}{\textwidth}{@{}lY Y@{}}
\toprule
Variable & Definition & Transformation and role \\
\midrule
Emergence this year & Indicator equal to one when the snapshot year is the emergence year for a city-genre cell & Binary outcome in exact-year fixed effects \\
Local spawning flow & Bands formed in a city-year whose founding cohort includes at least one musician with prior local band history in that city & Log transformed, one-year-lagged roll-5 average, z-scored; headline mechanism \\
Target-genre active bands & Active local bands in the focal broad genre family & Log transformed, one-year-lagged roll-5 average, z-scored; headline mechanism \\
Target-genre multi-band musicians & Local musicians active in the focal genre who hold at least two active local bands in that same genre & Log transformed, one-year-lagged roll-5 average, z-scored; headline mechanism \\
Overall active bands & Active local bands of any genre in the city-year snapshot & Log transformed, one-year-lagged roll-5 average, z-scored; control \\
Nontrivial communities & Number of detected communities with at least two active bands in the local yearly network & One-year-lagged roll-5 average, z-scored; control \\
Overall multi-band musicians & Active local musicians with at least two active local bands of any genre & Log transformed, one-year-lagged roll-5 average, z-scored; control \\
\bottomrule
\end{tabularx}

\vspace{0.4em}
\parbox{0.96\textwidth}{\footnotesize Notes: The preferred specification uses city-genre and year
fixed effects with standard errors clustered by city. All reported continuous predictors are
z-scored within the estimation sample.}
\end{table}


The preferred exact-year specification is a reduced-form transition equation for the probability
that a city-genre cell that has not yet emerged first crosses the paper's operational scene
threshold in year $t$:
\[
\Pr(E_{cgt}=1 \mid E_{cg,t-1}=0) = F\left(\alpha_{cg} + \lambda_t + \beta_1 S_{c,t-1}
 + \beta_2 B_{cg,t-1} + \beta_3 M_{cg,t-1}\right),
\]
where $E_{cgt}$ indicates whether city $c$ and genre $g$ emerge in year $t$, $S_{c,t-1}$ is
lagged local spawning flow, $B_{cg,t-1}$ is lagged target-genre active band stock, and
$M_{cg,t-1}$ is lagged target-genre multi-band musician depth. The design absorbs city-genre fixed
effects and year fixed effects, with standard errors clustered by city. This specification is
useful because it compares the same city-genre environment over time rather than relying on static
differences between famous and obscure scenes. It should therefore be read as a within-cell
transition equation on the risk set, not as a structural production function for scenes.

The design should not be read as delivering causal identification. Stronger local scenes may
generate some of the measured predictors, and omitted local shocks may move both current scene
conditions and later emergence. The coefficients therefore support a disciplined statement about
which local scene features predict later transitions into operational scene status, not a claim
that the paper directly observes knowledge transfer or the exact historical birth date of a genre.
